% Lösungen Einheit 1 von 4 – Die Schar als Familie (zusammenbau v0.1)
\einheitenkopf{Einheit 1 von 4 · Die Schar als Familie}

\erg{9}{a) Richtungsvektor – die Richtung ändert sich}
\erg{10}{a) für alle Mitglieder – $k$ muss beim Einsetzen herausfallen \quad b) $2k + 2 \cdot 3 - 2k = 6$ für jedes $k$ – $k$ fällt heraus, $P$ liegt in jeder $E_k$ \quad c) $E_5: x_1 + 5x_2 + x_3 = 10$; mit $x_1 = x_3 = 0$ gilt $5 \cdot 2 = 10$ \quad d) ja: aus $x_2$ Faktor $2$, dann $6 = 2k$ und $8 = 2(k + 1)$ mit demselben $k = 3$, rechte Seite $8 = 2 \cdot 4$ – $F$ ist $E_3$ \quad e) parallel und verschieden: der Richtungsvektor $(1 | 3 | 2)$ hängt nicht von $k$ ab; der Stützvektor ist $(0 | 1 | 0) + k \cdot (2 | 0 | -1)$, und $(2 | 0 | -1)$ ist kein Vielfaches von $(1 | 3 | 2)$ \quad f) $g: \vec{x} = (1 | 2 | 3) + k \cdot (1 | 0 | -2)$ \quad g) abgelesen $P(0 | 3 | 0)$ und $Q(0 | 0 | 1)$; $k \cdot 0 + 3 + 0 = 3$ und $k \cdot 0 + 0 + 3 = 3$ für jedes $k$ – wegen $x_1 = 0$ fällt $k$ heraus, beide Punkte und damit $h$ liegen in jeder $E_k$ \quad h) $x_1x_3$-Ebene – der Normalenvektor hat für jedes $k$ keine $x_2$-Komponente \quad i) Annahme $(a | 3 - a | 0) = r \cdot (b | 3 - b | 0)$ mit $a \ne b$: aus $a = rb$ und $3 - a = 3r - rb$ folgt $3 = 3r$, also $r = 1$ und $a = b$ – Widerspruch}
\erg{11}{a) Fehler: zwei Parameterwerte beweisen nichts für alle $k$. Richtig allgemein: $2k + 2 + 2(1 - k) = 4$ für jedes $k$ – der Parameter fällt heraus \quad b) Die entstehende Gleichung enthält den Parameter nicht mehr und ist wahr – also gilt sie für jeden Parameterwert, jedes Mitglied enthält den Punkt.}
